\section{需要多少资源，哪些地方值得压缩}\label{sec:resource}
\begin{table}[htbp]\centering\small
\caption{两层四存储的一轮Grover；不含真实准备、复位与读出}\label{tab:resources}
\begin{tabular}{@{}lrr@{}}\toprule
资源 & R5 & R6\\\midrule
总角色数 & 17 & 17\\
router调用 & 8 & 8\\
router内部CZ & 40 & 48\\
CZ总数 & 105 & 113\\
其中跨片CZ & 18 & 18\\
物理单门 & 288 & 272\\
CZ层深 & 86 & 92\\
总调度层（含入口等待） & 283 & 289\\
query时长/ns & 11184.152 & 11522.380\\
含扩散总时长/ns & 11426.152 & 11764.380\\
\bottomrule\end{tabular}\end{table}

R5/R6一轮分别105/113个CZ：节点部分是$8\times5/8\times6=40/48$，其余65个来自树控制保存与取出18、跨片完整SWAP 18、QPU中转SWAP 24、存储访问4、扩散1。跨片CZ已在总数中，不能重复加。

当前20个完整SWAP贡献60个CZ和120个半门。即使节点已经较短，接口与地址处理仍占很大比例。接收端已知0时可研究两CNOT搬运替换，但须重新证明全程入口、回程与相位，不能在每个位置直接删第三个CNOT。
\datafig{07_cost}{同一拓扑下的资源。10-CZ柱只是假定每次router核用10个CZ、保留其他协议的估算（$8\times10+65=145$），没有在此验证其脉冲性能，也不是最优门数证明。}{fig:cost}
对n层、$N=2^n$存储，本协议为
\begin{equation}
 N_{\rm role}=4N+n-1,\qquad N_{\rm router}=4N-2n-4.
\end{equation}
n为1、2、3时，角色数8/17/34，router调用2/8/22。这是串行地址注入、显式端口方案的实际计数，不是下界。
\datafig{08_resources_scale}{1至3层资源。时间统一比较query；CZ总数只有n=2含一次扩散，该点比较纯查询时应减1。H均已实际展开。}{fig:resourcescale}
\FloatBarrier

\section{数值可信度与结论边界}\label{sec:boundary}
\subsection{积分、能级与张量网络}
主局域步长0.1 ns，另固定参数用0.05 ns，并将四能级增加至五能级，没有重新优化。图\ref{fig:convergence}的变化小于百分比量级残差，支持局域传播与截断足够稳定。
\datafig{35_pulse_convergence}{四种上下行的积分与能级检查。五能级不是新的优化门，只是同一冻结波形的截断敏感性检验。}{fig:convergence}
整网使用bond 128，另对目标10用64重算；R5/R6查询$\Fe$差约$2.18\times10^{-6}$与$3.79\times10^{-5}$。同时检查正性、迹、辅助选择界与完整原生矩阵。累计截断权重只是诊断，不能直接当作保真度误差。

\subsection{扩大规模仍会累积相干残差}
固定均匀地址、D负号态，测试全1与单目标两种数据库模式。图\ref{fig:scaling}说明无耗散时模块相干偏差仍会累积。三层没有覆盖全部数据库、全部输入，也没有严格误差界。
\datafig{34_scaling}{R5的1至3层相位query。三层原始目标重叠约96.5\%--97.1\%，均用bond 128；只是两种数据库模式的一次query，不是多轮Grover或可扩展性证明。}{fig:scaling}

\subsection{真实器件还缺哪层证据}
最重要的缺口是共享工作点。根输出与子节点中心为同一物理位，不能同时赋予两个参考角色的不同频率。真实系统应按节点实际工作点分别校准，必要时加入转场、等待相位和帧转换。

QPU搬运、跨片、存储访问及H展开尚未全部换成真实脉冲；外部AD几乎无泄露受理想01延拓限制。当前只研究基态经典存储，没有证明任意量子数据库读写。

这里只验证一轮，不能将成功率简单乘方预测重复查询。小的末态辅助残留会改变下一次部分router入口，必须直接验证多次复用。
\FloatBarrier

\section{综合判断与实验推进顺序}\label{sec:conclusion}
线路结构上，R5/R6在规定入口可构造正确路由；完整协议保持这些入口条件，固定访问与相干恢复给出query，四目标Grover闭环成立。

参考控制上，局域四能级与整树模块过程接近目标，但仍有辅助恢复和相位残差。V配对处理当前辅助路径的目标旋转，端点帧处理剩余分支相位；两者不能替代泄露与路径错误的修复。

建议按因果顺序推进：先确认17角色拓扑、共享端口与各节点工作点；再校准01/12门、CZ及含2态搬运块；随后做各节点down、独立up和实际往返；再做空query、路径Ramsey与辅助恢复；最后固定数据库query、GHZ和Grover。

这个顺序先排人口和回程，再排相位，最后进入算法，使失败有较小定位范围。当前结论是具备值得落地验证的线路，而非真实17位器件已被证明达到本文百分数。
